% GENERATED FILE. Edit canonical lessons, then run npm run generate:books.
% canonical-source-hash: fnv1a64:438a032599a7f235
% canonical-lessons: HI-C319-pitna, HI-C319-jharap, HI-C319-ghusa, HI-C319-thappar, HI-C319-ainthna, HI-R319-review-312-319

\chapter{A Fight in the Market}
\label{ch:hi-a-fight-in-the-market}
\hlchaptermodality{319}

\begin{chapteropening}
\textbf{By the end of this chapter:} \emph{I can tell who hit whom in a fight.}

Complete the last lesson of chapter 319: I can tell who hit whom in a fight.
\end{chapteropening}

\section[\texorpdfstring{\hi{पीटना} (pīṭnā)}{पीटना (pīṭnā)}]{\hi{पीटना} (pīṭnā) — to beat}

\label{lesson:HI-C319-pitna}



\begin{quote}
Before the new one: say the Hindi for to quake, then the Hindi for a shock (of grief or fright).
\end{quote}

\subsection*{You'll want to know: \hi{पीटना}}
\textbf{\hi{पीटना}} — \emph{pīṭnā} — \textquotedblleft{}to beat\textquotedblright{}.

\begin{quote}
\hi{दो} \hi{आदमियों} \hi{ने} \hi{एक} \hi{चोर} \hi{को} \hi{पीटा।}
\end{quote}

\begin{quote}
\emph{do ādmiyoṁ ne ek cor ko pīṭā.}
\end{quote}

\begin{quote}
Two men beat up a thief.
\end{quote}

Past, with \hi{ने} (he or she …): \hi{उसने} \hi{ढोल} \hi{पीटा}, \emph{usne ḍhol pīṭā}.

\begin{cousinweb}
goes back to Sanskrit \emph{piṭṭayati}, to pound into a flat mass, beside which CDIAL reconstructs \emph{*piṭṭyatē}, \textquotedblleft{}is struck\textquotedblright{} (CDIAL 8165). It is hitting again and again; you also \textbf{\hi{पीटना}} a drum: \textbf{\hi{ढोल} \hi{पीटना}}.
\end{cousinweb}

\begin{grammarlens}[title={What you've built}]
One more word for telling what happened.
\end{grammarlens}

\subsection*{Your turn}
\begin{itemize}
\raggedright
  \item \emph{Say it:} \emph{pīṭnā}
  \item \emph{Say it:} \emph{pīṭnā}, once more
  \item \emph{Say it:} \emph{do ādmiyoṁ ne ek cor ko pīṭā.} — Two men beat up a thief.
  \item \emph{Recall:} say the Hindi for to quake, then the Hindi for a shock (of grief or fright), then say \emph{pīṭnā} again
\end{itemize}

\begin{tcolorbox}[breakable,colback=teal!4,colframe=teal!35!black,title={Before you move on}]
What does \hi{पीटना} mean? (\textquotedblleft{}To beat\textquotedblright{}.) Say it once more.
\end{tcolorbox}
\section[\texorpdfstring{\hi{झड़प} (jhaṛap)}{झड़प (jhaṛap)}]{\hi{झड़प} (jhaṛap) — a scuffle; a skirmish}

\label{lesson:HI-C319-jharap}



\begin{quote}
Before the new one: say the Hindi for a shock (of grief or fright), then the Hindi for to beat.
\end{quote}

\subsection*{You'll want to know: \hi{झड़प}}
\textbf{\hi{झड़प}} — \emph{jhaṛap} — \textquotedblleft{}a scuffle; a skirmish\textquotedblright{}.

\begin{quote}
\hi{बाज़ार} \hi{में} \hi{दो} \hi{दुकानदारों} \hi{के} \hi{बीच} \hi{झड़प} \hi{हुई।}
\end{quote}

\begin{quote}
\emph{bāzār meṁ do dukāndāroṁ ke bīc jhaṛap huī.}
\end{quote}

\begin{quote}
A scuffle broke out between two shopkeepers in the market.
\end{quote}

Feminine.

\begin{cousinweb}
is set by CDIAL under \emph{*jhaṭ}, a form it reconstructs for quick, sudden movements (CDIAL 5327), where it gives Hindi \emph{jhaṛap}, a fight: the same family as \textbf{\hi{झटका}}, a jolt. A \textbf{\hi{झड़प}} is short and sudden, between two people or two small groups.
\end{cousinweb}

\begin{grammarlens}[title={What you've built}]
One more word for telling what happened.
\end{grammarlens}

\subsection*{Your turn}
\begin{itemize}
\raggedright
  \item \emph{Say it:} \emph{jhaṛap}
  \item \emph{Say it:} \emph{jhaṛap}, once more
  \item \emph{Say it:} \emph{bāzār meṁ do dukāndāroṁ ke bīc jhaṛap huī.} — A scuffle broke out between two shopkeepers in the market.
  \item \emph{Recall:} say the Hindi for a shock (of grief or fright), then the Hindi for to beat, then say \emph{jhaṛap} again
\end{itemize}

\begin{tcolorbox}[breakable,colback=teal!4,colframe=teal!35!black,title={Before you move on}]
What does \hi{झड़प} mean? (\textquotedblleft{}A scuffle; a skirmish\textquotedblright{}.) Say it once more.
\end{tcolorbox}
\section[\texorpdfstring{\hi{घूँसा} (ghū̃sā)}{घूँसा (ghū̃sā)}]{\hi{घूँसा} (ghū̃sā) — a punch}

\label{lesson:HI-C319-ghusa}



\begin{quote}
Before the new one: say the Hindi for to beat, then the Hindi for a scuffle.
\end{quote}

\subsection*{You'll want to know: \hi{घूँसा}}
\textbf{\hi{घूँसा}} — \emph{ghū̃sā} — \textquotedblleft{}a punch\textquotedblright{}.

\begin{quote}
\hi{उसने} \hi{चोर} \hi{को} \hi{एक} \hi{घूँसा} \hi{मारा।}
\end{quote}

\begin{quote}
\emph{usne cor ko ek ghū̃sā mārā.}
\end{quote}

\begin{quote}
He punched the thief once.
\end{quote}

Masculine.

\begin{cousinweb}
\textbf{\hi{घूँसा}} (\emph{ghū̃sā}) is inherited from \emph{*ghūssa}, a form CDIAL reconstructs (CDIAL 4498). You \textbf{\hi{घूँसा} \hi{मारना}}, strike a punch: the punch is the object, so the verb is masculine.
\end{cousinweb}

\begin{grammarlens}[title={What you've built}]
One more word for telling what happened.
\end{grammarlens}

\subsection*{Your turn}
\begin{itemize}
\raggedright
  \item \emph{Say it:} \emph{ghū̃sā}
  \item \emph{Say it:} \emph{ghū̃sā}, once more
  \item \emph{Say it:} \emph{usne cor ko ek ghū̃sā mārā.} — He punched the thief once.
  \item \emph{Recall:} say the Hindi for to beat, then the Hindi for a scuffle, then say \emph{ghū̃sā} again
\end{itemize}

\begin{tcolorbox}[breakable,colback=teal!4,colframe=teal!35!black,title={Before you move on}]
What does \hi{घूँसा} mean? (\textquotedblleft{}A punch\textquotedblright{}.) Say it once more.
\end{tcolorbox}
\section[\texorpdfstring{\hi{थप्पड़} (thappaṛ)}{थप्पड़ (thappaṛ)}]{\hi{थप्पड़} (thappaṛ) — a slap}

\label{lesson:HI-C319-thappar}



\begin{quote}
Before the new one: say the Hindi for a scuffle, then the Hindi for a punch.
\end{quote}

\subsection*{You'll want to know: \hi{थप्पड़}}
\textbf{\hi{थप्पड़}} — \emph{thappaṛ} — \textquotedblleft{}a slap\textquotedblright{}.

\begin{quote}
\hi{मीरा} \hi{ने} \hi{जेबकतरे} \hi{को} \hi{थप्पड़} \hi{मारा।}
\end{quote}

\begin{quote}
\emph{Mīrā ne jebkatre ko thappaṛ mārā.}
\end{quote}

\begin{quote}
Mira slapped the pickpocket.
\end{quote}

Masculine.

\begin{cousinweb}
\textbf{\hi{थप्पड़}} (\emph{thappaṛ}) is inherited from \emph{*thapp}, a flat blow, a form CDIAL reconstructs (CDIAL 6091). Like \textbf{\hi{घूँसा}}, it is struck: \textbf{\hi{थप्पड़} \hi{मारा}}.
\end{cousinweb}

\begin{grammarlens}[title={What you've built}]
One more word for telling what happened.
\end{grammarlens}

\subsection*{Your turn}
\begin{itemize}
\raggedright
  \item \emph{Say it:} \emph{thappaṛ}
  \item \emph{Say it:} \emph{thappaṛ}, once more
  \item \emph{Say it:} \emph{Mīrā ne jebkatre ko thappaṛ mārā.} — Mira slapped the pickpocket.
  \item \emph{Recall:} say the Hindi for a scuffle, then the Hindi for a punch, then say \emph{thappaṛ} again
\end{itemize}

\begin{tcolorbox}[breakable,colback=teal!4,colframe=teal!35!black,title={Before you move on}]
What does \hi{थप्पड़} mean? (\textquotedblleft{}A slap\textquotedblright{}.) Say it once more.
\end{tcolorbox}
\section[\texorpdfstring{\hi{ऐंठना} (aiṇṭhnā)}{ऐंठना (aiṇṭhnā)}]{\hi{ऐंठना} (aiṇṭhnā) — to twist; to wring}

\label{lesson:HI-C319-ainthna}



\begin{quote}
Before the new one: say the Hindi for a punch, then the Hindi for a slap.
\end{quote}

\subsection*{You'll want to know: \hi{ऐंठना}}
\textbf{\hi{ऐंठना}} — \emph{aiṇṭhnā} — \textquotedblleft{}to twist; to wring\textquotedblright{}.

\begin{quote}
\hi{सिपाही} \hi{ने} \hi{चोर} \hi{का} \hi{हाथ} \hi{ऐंठा।}
\end{quote}

\begin{quote}
\emph{sipāhī ne cor kā hāth aiṇṭhā.}
\end{quote}

\begin{quote}
The constable twisted the thief's arm.
\end{quote}

Past, with \hi{ने} (he or she …): \hi{उसने} \hi{रस्सी} \hi{ऐंठी}, \emph{usne rassī aiṇṭhī}.

\begin{cousinweb}
\textbf{\hi{ऐंठना}} (\emph{aiṇṭhnā}) is set by CDIAL under Sanskrit \emph{āveṣṭate}, of the family of \emph{veṣṭ}, to wind round (CDIAL 1448), where Hindi is \emph{aĩṭhnā}. It begins with \textbf{\hi{ऐ}}, the vowel of \textbf{\hi{ऐसा}}. You \textbf{\hi{ऐंठते} \hi{हो}} a rope, an arm or a wet cloth.
\end{cousinweb}

\begin{grammarlens}[title={What you've built}]
One more word for telling what happened.
\end{grammarlens}

\subsection*{Your turn}
\begin{itemize}
\raggedright
  \item \emph{Say it:} \emph{aiṇṭhnā}
  \item \emph{Say it:} \emph{aiṇṭhnā}, once more
  \item \emph{Say it:} \emph{sipāhī ne cor kā hāth aiṇṭhā.} — The constable twisted the thief's arm.
  \item \emph{Recall:} say the Hindi for a punch, then the Hindi for a slap, then say \emph{aiṇṭhnā} again
\end{itemize}

\begin{tcolorbox}[breakable,colback=teal!4,colframe=teal!35!black,title={Before you move on}]
What does \hi{ऐंठना} mean? (\textquotedblleft{}To twist; to wring\textquotedblright{}.) Say it once more.
\end{tcolorbox}
\section[(dialogue)]{Looking back: from The Morning News to A Fight in the Market}

\label{lesson:HI-R319-review-312-319}



\begin{quote}
Say the Hindi for a radio and the news (formal).
\end{quote}

\subsection*{Your turn}
Say these aloud:

\begin{itemize}
\raggedright
  \item a radio, the news (formal), a broadcast, a headline, to be printed
  \item a journalist, the scene (of an incident), to print, an editor, a detailed account
  \item a rumour, talk, terror, to stir up, a sensation
  \item a claim, the public, a declaration, information, to put off
  \item a disturbance, trouble, to burst, a staff, a curfew
  \item a mystery, an informer, a guess, to be startled, astonished
  \item a bang, a bomb, to echo, to quake, a shock (of grief or fright)
  \item to beat, a scuffle, a punch, a slap, to twist
\end{itemize}

\begin{tcolorbox}[breakable,colback=teal!4,colframe=teal!35!black,title={Before you move on}]
A radio? (\textbf{\hi{रेडियो}}, \emph{reḍiyo}.) The news (formal)? (\textbf{\hi{समाचार}}, \emph{samācār}.) A broadcast? (\textbf{\hi{प्रसारण}}, \emph{prasāraṇ}.) A headline? (\textbf{\hi{सुर्ख़ी}}, \emph{surkhī}.) To be printed? (\textbf{\hi{छपना}}, \emph{chapnā}.) A journalist? (\textbf{\hi{पत्रकार}}, \emph{patrakār}.) The scene (of an incident)? (\textbf{\hi{घटनास्थल}}, \emph{ghaṭnāsthal}.) To print? (\textbf{\hi{छापना}}, \emph{chāpnā}.) An editor? (\textbf{\hi{संपादक}}, \emph{sampādak}.) A detailed account? (\textbf{\hi{विवरण}}, \emph{vivaraṇ}.) A rumour? (\textbf{\hi{अफ़वाह}}, \emph{afvāh}.) Talk? (\textbf{\hi{चर्चा}}, \emph{carcā}.) Terror? (\textbf{\hi{दहशत}}, \emph{dahśat}.) To stir up? (\textbf{\hi{भड़काना}}, \emph{bhaṛkānā}.) A sensation? (\textbf{\hi{सनसनी}}, \emph{sansanī}.) A claim? (\textbf{\hi{दावा}}, \emph{dāvā}.) The public? (\textbf{\hi{जनता}}, \emph{jantā}.) A declaration? (\textbf{\hi{घोषणा}}, \emph{ghoṣṇā}.) Information? (\textbf{\hi{जानकारी}}, \emph{jānkārī}.) To put off? (\textbf{\hi{टालना}}, \emph{ṭālnā}.) A disturbance? (\textbf{\hi{उपद्रव}}, \emph{upadrav}.) Trouble? (\textbf{\hi{फ़साद}}, \emph{fasād}.) To burst? (\textbf{\hi{फूटना}}, \emph{phūṭnā}.) A staff? (\textbf{\hi{लाठी}}, \emph{lāṭhī}.) A curfew? (\textbf{\hi{कर्फ़्यू}}, \emph{karfyū}.) A mystery? (\textbf{\hi{रहस्य}}, \emph{rahasya}.) An informer? (\textbf{\hi{मुखबिर}}, \emph{mukhbir}.) A guess? (\textbf{\hi{अंदाज़ा}}, \emph{andāzā}.) To be startled? (\textbf{\hi{चौंकना}}, \emph{cauṅknā}.) Astonished? (\textbf{\hi{हैरान}}, \emph{hairān}.) A bang? (\textbf{\hi{धमाका}}, \emph{dhamākā}.) A bomb? (\textbf{\hi{बम}}, \emph{bam}.) To echo? (\textbf{\hi{गूँजना}}, \emph{gū̃jnā}.) To quake? (\textbf{\hi{थर्राना}}, \emph{tharrānā}.) A shock (of grief or fright)? (\textbf{\hi{सदमा}}, \emph{sadmā}.) To beat? (\textbf{\hi{पीटना}}, \emph{pīṭnā}.) A scuffle? (\textbf{\hi{झड़प}}, \emph{jhaṛap}.) A punch? (\textbf{\hi{घूँसा}}, \emph{ghū̃sā}.) A slap? (\textbf{\hi{थप्पड़}}, \emph{thappaṛ}.) To twist? (\textbf{\hi{ऐंठना}}, \emph{aiṇṭhnā}.)
\end{tcolorbox}
